\chapter{Cranelift Native Code Generation}
\label{chap:cranelift}

\section{The Cranelift Backend Architecture}
\label{sec:cranelift:arch}

NumLang employs Cranelift \citep{cranelift2024} as its default native code generation engine (\texttt{src/codegen/cranelift/}). Cranelift is an optimizing, value-based compiler designed specifically for high-throughput code generation, memory safety, and sub-millisecond cold start times.

The Cranelift backend comprises six core subsystems:
\begin{itemize}
    \item \texttt{mod.rs}: Backend entry point, target machine configuration, and compilation orchestration.
    \item \texttt{abi.rs}: System V and Windows x86-64 calling convention mappings, struct layouts, and stack slot management.
    \item \texttt{mir\_emit.rs}: SSA MIR to Cranelift Intermediate Language (CLIF) lowering.
    \item \texttt{intrinsics.rs}: Mathematical routines, syscalls, and recurrence companion drivers.
    \item \texttt{deopt.rs}: Speculative deoptimization trampolines and register mapping tables.
    \item \texttt{linker.rs}: Native object file emission and platform linker orchestration.
\end{itemize}

\section{Strict Error Discipline: Zero Panics in Code Generation}
\label{sec:cranelift:zero_panic}

In accordance with NumLang's foundational computational integrity rules, the entire code generation pipeline enforces a strict invariant: \textbf{zero \texttt{panic!()}, zero \texttt{.unwrap()}, and zero \texttt{.expect()} calls}.

Every compilation failure---ranging from unsupported platform target features to malformed SSA basic block signatures---is returned as a structured, localized \texttt{CodegenError} variant:

\begin{lstlisting}[language=Rust, caption={The CodegenError enumeration in NumLang.}]
pub enum CodegenError {
    BackendError(String),
    TypeMismatch { expected: String, found: String },
    UnknownVariable(String),
    UnknownFunction(String),
    InvalidTerminator(String),
    UnsupportedFeature(String),
    LinkingFailed(String),
}
\end{lstlisting}

\section{SSA MIR to CLIF Translation (mir\_emit.rs)}
\label{sec:cranelift:mir_emit}

The lowering engine (\texttt{src/codegen/cranelift/mir\_emit.rs}) translates SSA MIR basic blocks into CLIF functions in two linear passes:

\subsection{Pass 1: Block and Phi Declaration}
Before translating instructions, the generator scans all MIR basic blocks. For each basic block, it creates a corresponding Cranelift \texttt{Block}. Any SSA $\phi$-nodes are declared as block parameters, guaranteeing that phi-node value cycles are resolved cleanly without introducing temporary stack spills.

\subsection{Pass 2: Instruction Selection and Branch Emission}
Instructions are emitted sequentially within their respective blocks:
\begin{itemize}
    \item \textbf{Binary Operations}: Arithmetic operations (\texttt{add}, \texttt{sub}, \texttt{mul}) map directly to native CLIF 64-bit integer instructions (\texttt{iadd}, \texttt{isub}, \texttt{imul}).
    \item \textbf{Recurrence Intrinsics}: Residual recurrence companion calls lower to high-speed binary exponentiation loops emitted directly into native machine code.
    \item \textbf{Terminators}: \texttt{Branch} lowers to \texttt{ins().jump()}; \texttt{BranchIf} lowers to \texttt{ins().brnz()}; \texttt{Return} lowers to \texttt{ins().return\_()}.
\end{itemize}

\section{The Scoped Arena Allocator Runtime}
\label{sec:cranelift:arena}

To support lightning-fast heap allocations during recursive loops, NumLang integrates a native scoped arena allocator (\texttt{src/runtime/arena.rs}, \texttt{src/runtime/arena.c}).

When the supercompiler determines that a recursive knot does not allow allocated objects to escape beyond the loop lifetime, it wraps the loop in an arena scope:
\begin{enumerate}
    \item Heap allocations invoke \texttt{\_\_numlang\_alloc}, which increments a thread-local bump pointer in $\mathcal{O}(1)$ time with zero mutex contention.
    \item Upon loop termination, \texttt{\_\_numlang\_arena\_reset} restores the pointer to the base watermark, reclaiming all loop-allocated memory in a single instruction without garbage collection overhead.
\end{enumerate}

\section{CLIF Instruction Lowering Specifications}
\label{sec:cranelift:clif_specs}

Table~\ref{tab:clif_lowering} outlines the exact semantic mappings from NumLang SSA MIR instructions to Cranelift CLIF primitives implemented in \texttt{src/codegen/cranelift/mir\_emit.rs}.

\begin{table}[h!]
\centering
\small
\begin{tabular}{lll}
\toprule
\textbf{SSA MIR Instruction} & \textbf{Cranelift CLIF Instruction} & \textbf{Semantic Description} \\
\midrule
\texttt{Constant(c)} & \texttt{ins().iconst(I64, c)} & 64-bit integer immediate \\
\texttt{BinOp(Add, x, y)} & \texttt{ins().iadd(vx, vy)} & Wrapping integer addition \\
\texttt{BinOp(Sub, x, y)} & \texttt{ins().isub(vx, vy)} & Wrapping integer subtraction \\
\texttt{BinOp(Mul, x, y)} & \texttt{ins().imul(vx, vy)} & Wrapping integer multiplication \\
\texttt{BinOp(Div, x, y)} & \texttt{ins().sdiv(vx, vy)} & Signed integer division \\
\texttt{BinOp(BitAnd, x, y)} & \texttt{ins().band(vx, vy)} & Bitwise logical AND \\
\texttt{BinOp(BitOr, x, y)} & \texttt{ins().bor(vx, vy)} & Bitwise logical OR \\
\texttt{BinOp(BitXor, x, y)} & \texttt{ins().bxor(vx, vy)} & Bitwise logical XOR \\
\texttt{Alloc(sz)} & \texttt{ins().call(malloc, [sz])} & Heap allocation \\
\texttt{Load(ptr)} & \texttt{ins().load(I64, flags, ptr, 0)} & 64-bit memory dereference \\
\texttt{Store(ptr, val)} & \texttt{ins().store(flags, val, ptr, 0)} & 64-bit memory write \\
\texttt{Branch(bb)} & \texttt{ins().jump(target\_block, [])} & Direct unconditional jump \\
\texttt{BranchIf(c, t, f)} & \texttt{ins().brnz(vc, target\_t, [])} & Conditional branch \\
\texttt{Return(val)} & \texttt{ins().return\_([val])} & Return from function \\
\bottomrule
\end{tabular}
\caption{Exact translation mappings from NumLang SSA MIR to Cranelift CLIF.}
\label{tab:clif_lowering}
\end{table}

\section{The Native Scoped Arena Implementation (arena.rs)}
\label{sec:cranelift:arena_code}

Listing~\ref{lst:arena_code} presents the scoped bump allocator runtime from \texttt{src/runtime/arena.rs}:

\begin{lstlisting}[language=Rust, caption={Scoped Arena Bump Allocator Runtime (\texttt{src/runtime/arena.rs}).}, label={lst:arena_code}]
pub const DEFAULT_CHUNK_CAPACITY: usize = 64 * 1024; // 64 KiB default chunk size

#[repr(C)]
pub struct ArenaChunk {
    pub data: *mut u8,
    pub capacity: usize,
    pub offset: usize,
    pub next: *mut ArenaChunk,
}

#[repr(C)]
pub struct Arena {
    pub head: *mut ArenaChunk,
    pub current: *mut ArenaChunk,
    pub default_chunk_size: usize,
    pub total_allocated: usize,
    pub peak_usage: usize,
}

impl Arena {
    pub fn new(capacity: usize) -> Self {
        let cap = if capacity == 0 { DEFAULT_CHUNK_CAPACITY } else { capacity };
        let chunk = Self::alloc_chunk(cap);
        Arena {
            head: chunk,
            current: chunk,
            default_chunk_size: cap,
            total_allocated: 0,
            peak_usage: 0,
        }
    }

    fn alloc_chunk(capacity: usize) -> *mut ArenaChunk {
        unsafe {
            let data_layout = Layout::from_size_align_unchecked(capacity, 16);
            let data = alloc(data_layout);
            assert!(!data.is_null(), "Arena: out of memory allocating chunk data");

            let chunk_layout = Layout::new::<ArenaChunk>();
            let chunk_ptr = alloc(chunk_layout) as *mut ArenaChunk;
            assert!(!chunk_ptr.is_null(), "Arena: out of memory allocating chunk header");

            ptr::write(chunk_ptr, ArenaChunk {
                data,
                capacity,
                offset: 0,
                next: ptr::null_mut(),
            });
            chunk_ptr
        }
    }

    pub fn alloc(&mut self, size: usize) -> *mut u8 {
        if size == 0 {
            return ptr::null_mut();
        }
        let aligned_size = (size + 7) & !7; // 8-byte alignment

        unsafe {
            let curr = self.current;
            if (*curr).offset + aligned_size <= (*curr).capacity {
                let res = (*curr).data.add((*curr).offset);
                (*curr).offset += aligned_size;
                self.total_allocated += aligned_size;
                if self.total_allocated > self.peak_usage {
                    self.peak_usage = self.total_allocated;
                }
                return res;
            }

            // Need next chunk or allocate a new one
            if !(*curr).next.is_null() && (*(*curr).next).capacity >= aligned_size {
                self.current = (*curr).next;
                let c = self.current;
                (*c).offset = aligned_size;
                self.total_allocated += aligned_size;
                if self.total_allocated > self.peak_usage {
                    self.peak_usage = self.total_allocated;
                }
                return (*c).data;
            }

            let new_cap = self.default_chunk_size.max(aligned_size * 2);
            let new_chunk = Self::alloc_chunk(new_cap);
            (*curr).next = new_chunk;
            self.current = new_chunk;
            (*new_chunk).offset = aligned_size;
            self.total_allocated += aligned_size;
            if self.total_allocated > self.peak_usage {
                self.peak_usage = self.total_allocated;
            }
            (*new_chunk).data
        }
    }

    pub fn reset(&mut self) {
        unsafe {
            let mut curr = self.head;
            while !curr.is_null() {
                (*curr).offset = 0;
                curr = (*curr).next;
            }
            self.current = self.head;
            self.total_allocated = 0;
        }
    }

    pub fn destroy(&mut self) {
        unsafe {
            let mut curr = self.head;
            while !curr.is_null() {
                let next = (*curr).next;
                let data_layout = Layout::from_size_align_unchecked((*curr).capacity, 16);
                dealloc((*curr).data, data_layout);
                let chunk_layout = Layout::new::<ArenaChunk>();
                dealloc(curr as *mut u8, chunk_layout);
                curr = next;
            }
            self.head = ptr::null_mut();
            self.current = ptr::null_mut();
            self.total_allocated = 0;
        }
    }
}

impl Drop for Arena {
    fn drop(&mut self) {
        self.destroy();
    }
}

\end{lstlisting}
